\documentclass[10pt]{article}
\usepackage{tmlr}
\usepackage{amsmath,amssymb}
\usepackage[T1]{fontenc}
\usepackage[utf8]{inputenc}
\usepackage{lmodern}
\usepackage{microtype}
\usepackage{booktabs}
\usepackage{xurl}
\usepackage[hidelinks]{hyperref}
\DeclareUnicodeCharacter{220E}{\ensuremath{\square}}
\setlength{\emergencystretch}{3em}
\title{When a Representation Can Certify: Sharp Fibre-Diameter Limits and Minimal Refinement}
\author{Anonymous authors}
\def\month{MM}
\def\year{YYYY}
\def\openreview{\url{https://openreview.net/forum?id=XXXX}}
\date{}
\begin{document}
\maketitle

\begin{abstract}
A representation can be sufficient for making a decision while remaining
insufficient for certifying the value attached to that decision. We give
an exact finite-fibre characterization of this gap. Let \(phi\) map
instances to a representation, let \(F_z\) be a representation fibre,
let \texttt{V} be the target quantity, and let \(D(z)\) be the range of
\texttt{V} on that fibre. Any deterministic point certificate that is
constant on \(F_z\) has worst-case error at least \texttt{D(z)/2}, and
any interval of radius below \texttt{D(z)/2} must miss at least one
fibre member. The bound is sharp.

The same quantity gives a constructive converse. A fibre admits an
\(eps\)-valid constant certificate if and only if \(D(z) \le 2 eps\);
the midpoint of its extreme target values is then valid. For finite
fibres, the minimum number of refined parts required for \(eps\)-valid
certification is exactly the count returned by a greedy interval sweep
over the sorted target values. Under a restricted separator family,
certification is possible if and only if the separators distinguish
every pair whose target values differ by more than \(2 eps\). Without
refinement, achievable coverage is exactly the mass of fibres already
satisfying \(D(z) \le 2 eps\). These results turn a
representation-induced impossibility boundary into an exact
refine-or-abstain law.

Exhaustive independent-code-path checks within this repository found no
violations over the specified finite configuration families and include
planted counterexamples to show that the checkers can fail. We also
retain a prospective empirical failure that motivated the theory: a
frozen certificate study reached \texttt{44/44} coverage but incurred
\texttt{20/44} strict held-out violations, and its available selection
score was not detectably correlated with realised excess (\(r=-0.144\),
permutation \(p=0.353\), \(n=44\)). Those data do not establish the
fibre diameters of the real corpus; they show why marginal coverage and
a weak selector cannot substitute for conditional certifiability. The
contribution is therefore a sharp representation-level certification
criterion and exact refinement calculus, not a claim of broad empirical
transfer.
\end{abstract}

\noindent\textbf{Keywords:} representation sufficiency; certification; fibre diameter; refinement; selective uncertainty

\hypertarget{introduction}{%
\section{1. Introduction}\label{introduction}}

Many learned and rule-based systems compress a complex object into a
representation before making a prediction, selecting an action, or
issuing a certificate. The representation may be excellent for one query
and inadequate for another. A state summary can identify which action to
take while failing to identify how much value that action carries. A
classifier can have good marginal performance while a conditional
certificate fails on the subset where it is actually used. A coarse
state can support a ranking but not a guarantee.

This paper asks a narrower question than general representation
learning: \textbf{when can a fixed representation support a uniformly
valid certificate, and what is the minimum refinement required when it
cannot?} The answer is governed by one object, the target diameter of
each representation fibre. If two instances are indistinguishable to the
certificate but their target values are far apart, no tuning of the
certificate can make that ambiguity disappear. Conversely, once
refinement reduces every fibre diameter below the target tolerance, a
valid certificate exists by construction.

This deterministic viewpoint complements statistical work on marginal,
conditional, and selection-conditional uncertainty. Conformal prediction
supplies distribution-free marginal guarantees under exchangeability,
while recent selective-conformal work studies how coverage changes after
data-dependent selection. Those methods address sampling validity. Our
question is prior to sampling: if a representation maps two
target-separated instances to the same observable state, then no
fibre-constant certificate can distinguish them, regardless of
calibration procedure. We therefore treat conditional coverage methods
as donors for the statistical layer and isolate the representation-level
obstruction underneath it.

The paper makes five claims. First, \texttt{D(z)/2} is the exact
worst-case radius of any deterministic point certificate accepted on a
fibre. Second, \(D(z) \le 2 eps\) is necessary and sufficient for an
\(eps\)-valid fibre-constant certificate. Third, the minimum
unconstrained refinement cost is exactly computable by a greedy interval
sweep. Fourth, a restricted separator family can realize certification
exactly when it separates every pair with target gap greater than
\(2 eps\). Fifth, without refinement the maximum certifiable coverage is
exactly the probability mass of already-small fibres, giving a
refine-or-abstain frontier.

These claims do \textbf{not} rely on the earlier \texttt{A\_t/B\_t}
gadget family from the previous manuscript. An adversarial proof review
found that those all-\texttt{t} formulas require an unstated
cross-gadget separability argument under a non-additive global width
term. V3 therefore removes that family, its minimax corollaries, the
conditional four-index compiler theorem, and the missing single-block
convention from the submission spine. They remain historical research
records, not premises of the present paper.

\hypertarget{setting-and-terminology}{%
\section{2. Setting and terminology}\label{setting-and-terminology}}

Let \texttt{X} be a finite instance set and let

\[phi : X -> Z\]

be a representation. For an observed state \texttt{z}, define the fibre

\(F_z = {x in X : phi(x) = z}\).

Let

\texttt{V\ :\ X\ -\textgreater{}\ R}

be the scalar target to be certified. Its fibre diameter is

\(D_phi(z) = max_{x,x' in F_z} |V(x)-V(x')|\).

A deterministic point certificate based only on \(phi\) is a function
\texttt{c\ :\ Z\ -\textgreater{}\ R}. It is accepted on a fibre when it
is issued for every member of that fibre. Because all members share the
same representation, the certificate must return the same value for all
of them.

An interval certificate is likewise a function of \texttt{z}; write its
centre as \(c(z)\) and radius as \(r(z)\). At a fixed tolerance
\(eps \ge 0\), a point certificate is \(eps\)-valid on \(F_z\) when

\[|c(z)-V(x)| \le  eps\]

for every \(x in F_z\).

A representation \(phi'\) refines \(phi\) when equality under \(phi'\)
implies equality under \(phi\). Equivalently, every \(phi'\) fibre is
contained in one \(phi\) fibre.

The theory is deliberately finite and deterministic. It assumes no
sampling model, prior, exchangeability, smoothness, learned model class,
or computational hardness. Statistical estimation of \(D(z)\) and
learning a useful refinement are separate problems.

\hypertarget{the-fibre-diameter-floor}{%
\section{3. The fibre-diameter floor}\label{the-fibre-diameter-floor}}

\hypertarget{point-certificates}{%
\subsection{3.1 Point certificates}\label{point-certificates}}

\textbf{Theorem 1 (sharp point-certificate floor).} For every fibre
\(F_z\) and every deterministic point certificate \(c(z)\) accepted on
it,

\texttt{max\_\{x\ in\ F\_z\}\ \textbar{}c(z)-V(x)\textbar{}\ \textgreater{}=\ D\_phi(z)/2}.

\textbf{Proof.} Choose two fibre members with target separation
\(D_phi(z)\). By the triangle inequality,

\(D_phi(z) \le |V(x)-c(z)| + |c(z)-V(x')|\).

At least one term is therefore at least \texttt{D\_phi(z)/2}.
Conversely, the midpoint

\texttt{c*(z)\ =\ (min\_\{x\ in\ F\_z\}\ V(x)\ +\ max\_\{x\ in\ F\_z\}\ V(x))/2}

has worst-case error exactly \texttt{D\_phi(z)/2}. The bound is exact. ∎

This is an information statement rather than a complexity statement.
Unlimited compute cannot beat the radius while the certificate sees only
\texttt{z}.

\hypertarget{interval-certificates}{%
\subsection{3.2 Interval certificates}\label{interval-certificates}}

\textbf{Theorem 2 (interval floor).} An interval of radius
\texttt{r\ \textless{}\ D\_phi(z)/2} cannot contain \(V(x)\) for every
\(x in F_z\).

The result follows because any interval covering all target values on
the fibre must span a set of diameter \(D_phi(z)\) and therefore have
width at least \(D_phi(z)\).

A probabilistic corollary follows immediately. If a conditional
distribution places probability one half on each endpoint of a
diameter-attaining pair, then every fibre-constant interval with radius
below \texttt{D\_phi(z)/2} has conditional miscoverage at least one
half. This is a worst-case witness, not a statement that every real
fibre has such a distribution.

\hypertarget{from-impossibility-to-construction}{%
\section{4. From impossibility to
construction}\label{from-impossibility-to-construction}}

The floor becomes useful when read as an exact design criterion rather
than only as a negative result.

\hypertarget{exact-certifiability-threshold}{%
\subsection{4.1 Exact certifiability
threshold}\label{exact-certifiability-threshold}}

\textbf{Theorem 3 (certifiability equivalence).} An \(eps\)-valid
deterministic point certificate constant on a fibre exists if and only
if

\(D_phi(z) \le 2 eps\).

Necessity is Theorem 1. Sufficiency is constructive: the midpoint
certificate lies within half the diameter of every fibre member.

Thus a representation is not ``approximately certifiable'' in an
informal sense. At a declared tolerance, each fibre is either already
certifiable or must be split.

\hypertarget{minimum-unconstrained-refinement}{%
\subsection{4.2 Minimum unconstrained
refinement}\label{minimum-unconstrained-refinement}}

Fix one fibre and sort its target values:

\(v_1 \le v_2 \le ... \le v_n\).

A refined part is certifiable at tolerance \(eps\) exactly when its
diameter is at most \(2 eps\). The problem is therefore to cover the
sorted values with the fewest intervals of length \(2 eps\).

\textbf{Theorem 4 (exact minimum refinement cost).} The minimum number
of certifiable parts equals the count returned by the following greedy
sweep: start a part at the smallest uncovered value and include every
subsequent value no more than \(2 eps\) above that part's minimum; then
repeat on the remaining suffix.

\textbf{Proof sketch.} An optimal part containing the smallest remaining
value cannot extend beyond that value plus \(2 eps\). The greedy first
part covers every point any feasible first part could cover. Replacing
an optimal first part by the greedy one cannot increase the number of
parts. Induction on the remaining suffix proves optimality. ∎

Write this minimum as \(k*(z,eps)\). The number \texttt{k*-1} is the
exact unconstrained number of additional representation states required
to make that fibre certifiable.

\hypertarget{what-available-separators-can-realize}{%
\subsection{4.3 What available separators can
realize}\label{what-available-separators-can-realize}}

An abstract partition may not be obtainable from the information a
system can actually compute. Let \texttt{S} be a family of admissible
predicates or features, and require the refined representation to be
measurable with respect to \texttt{S}.

Two instances are \texttt{S}-indistinguishable when every predicate in
\texttt{S} takes the same value on them.

\textbf{Theorem 5 (separator realizability).} An \texttt{S}-measurable
\(eps\)-valid refinement exists if and only if every
\texttt{S}-indistinguishable pair satisfies

\(|V(x)-V(x')| \le 2 eps\).

\textbf{Proof.} If an indistinguishable pair exceeds \(2 eps\), every
\texttt{S}-measurable representation places it in one atom, which is
uncertifiable by Theorem 3. Conversely, if no such pair exists, each
atom of the joint \texttt{S}-signature has target diameter at most
\(2 eps\); the midpoint certificate is valid on every atom. ∎

This theorem separates two costs that are easy to conflate. \texttt{k*}
is the information-theoretic minimum over arbitrary refinements. The
realizable cost depends on the separator vocabulary. A weak vocabulary
may require more states than \texttt{k*}, and if it merges one
target-separated pair it cannot certify the fibre at any number of
downstream tuning steps.

\hypertarget{the-refine-or-abstain-frontier}{%
\section{5. The refine-or-abstain
frontier}\label{the-refine-or-abstain-frontier}}

Suppose fibres have probability mass \(P(phi=z)\) under a target
population. Without refinement, a fibre can be certified at tolerance
\(eps\) exactly when \(D_phi(z) \le 2 eps\).

\textbf{Theorem 6 (coverage identity).} The maximum coverage achievable
by a certificate that either accepts an entire original fibre or
abstains on it is

\(sum_z P(phi=z) * 1{D_phi(z) \le 2 eps}\).

Refining an uncertifiable fibre purchases exactly that fibre's
probability mass at the representation cost required to split it into
certifiable parts. Under unconstrained refinement that cost is
\(k*(z,eps)-1\); under a separator family, it is the corresponding
measurable refinement cost.

This gives an exact frontier rather than a generic recommendation to
``use more features.'' For each fibre, the relevant questions are: how
large is the target diameter, which target-separated pairs the available
information can distinguish, and whether the purchased coverage is worth
the required refinement.

\hypertarget{finite-hostile-verification}{%
\section{6. Finite hostile
verification}\label{finite-hostile-verification}}

The mathematical statements are proved above. Separate exhaustive
in-repository checkers were used as hostile finite-model verification
rather than as theorem authority. These independent-code-path checks are
not external replication.

For the fibre-diameter floor, the specified checker enumerated 784
finite configurations. It found zero point certificates beating
\texttt{D/2}, zero intervals of radius below \texttt{D/2} covering both
diameter endpoints, and zero balanced two-point examples with
miscoverage below one half. Its negative controls give the certificate
access to the member identity and require the checker to detect that the
floor can then be beaten; planted violations fired rather than passing
vacuously.

For the refinement theorem, the specified checker covered 4,704 main
configurations plus nested separator enumerations. It compared the
greedy part count against an exhaustive enumeration of set partitions on
an independent code path. The specified checker claims had zero
violations. Planted-violation controls fired for the sufficiency,
necessity, separator, and coverage predicates, while the no-alarm
control stayed silent.

These finite enumerations test implementation and transcription. General
authority comes from the proofs, not from the absence of a small
counterexample.

\hypertarget{a-preserved-adverse-empirical-boundary}{%
\section{7. A preserved adverse empirical
boundary}\label{a-preserved-adverse-empirical-boundary}}

The theory was motivated by a sequence of failed certificate
constructions, and those failures remain part of the scientific record
rather than being overwritten by the later theorem.

The first counted held-out revival study asked whether a Hamming-radius
backoff could restore certified coverage on 44 held-out PMLB decisions.
Against the corrected exact-cell parent method at
\texttt{0/44\ =\ 0.0000}, the backoff reached \texttt{32/44\ =\ 0.7273}.
That improved coverage substantially but missed the frozen \texttt{0.95}
gate. The outcome-independent lexical negative control, which ignores
the geometry entirely, reached \texttt{39/44\ =\ 0.8864} --- a higher
raw count than the specified Hamming geometry. That gap is not
statistically established. The two arms are evaluated on the same 44
decisions under the same fold assignment, so the comparison is paired:
of the 13 datasets on which the arms disagree, the control certifies 10
and the geometry 3, giving an exact McNemar two-sided \(p = 0.0923\) and
a 20,000-replicate paired bootstrap interval on the difference of
\texttt{{[}-0.3182,\ 0.0000{]}}, which does not exclude zero. The
proposed geometry was therefore not validated by this round, but the
reason is that it missed the frozen \texttt{0.95} gate at
\texttt{0.7273}, not that the control outperformed it: the data do not
support the stronger claim that a geometry-free baseline beats the
geometry. Within the same round, the primary learned ordering minus the
matched static adaptive arm had mean excess difference
\texttt{-0.001218244987} with bootstrap 95\% interval
\texttt{{[}-0.011716227308,\ 0.008821064426{]}}, while acquiring more
groups on average (\texttt{1.113636} versus \texttt{0.840909}); it did
not establish learned-ordering value, and it carried 24 strict
realized-bound violations among 42 certified commits. The first study
therefore remained below its pre-specified coverage gate.

The second counted held-out certificate study evaluated the same 44
held-out decisions. Its arm-conditional construction raised coverage
from the first study's \texttt{32/44} to \texttt{44/44}, meeting the
pre-specified \texttt{0.95} gate --- and validity still failed: strict
held-out violations were \texttt{20/44\ =\ 0.455}, against a
pre-specified maximum of \texttt{0.10}. A matched no-geometry lexical
control also reached \texttt{44/44} coverage, with fewer violations than
the geometric arm (14 versus 20), so the specified geometry again
supplied no measured advantage on that corpus. The second study
therefore failed its pre-specified validity criterion.

Across both counted attempts the same outcome-independent negative
control was at least as strong as the specified geometry on the recorded
criteria. That repetition, rather than either round alone, is what the
theory section treats as the empirical boundary worth explaining.

A subsequent diagnostic on the same frozen records showed that the
pre-specified target was arithmetically infeasible at full coverage:
\texttt{11/44\ =\ 0.25} held-out excesses exceeded \(tau=0.02\), whereas
the gate required at most \texttt{0.10} violations with a bound no
larger than \texttt{tau}. A split-conformal bound valid at
\(alpha=0.10\) on all 44 points was \texttt{0.061381}, about
\texttt{3.07\ tau}.

An oracle abstention analysis showed that a useful interior point
existed: removing the 25\% highest realised-excess cases left 33 cases
with zero violations and a conformal bound \texttt{0.016837}, below
\texttt{tau}. This oracle is not an implementable method because
realised excess is unknown at decision time.

The available model score did not provide a substitute. Its Pearson
correlation with realised excess was \texttt{-0.1442} (\(p=0.3528\)
under 20,000 permutations; Spearman \(rho=-0.1921\), \(n=44\)).
Abstaining on that score made the retained violation rate worse. A
noise-degraded oracle simulation suggested that correlations near
\texttt{0.85} would be needed on that one empirical distribution, but
the value is a design target for a successor, not a general constant.

These observations are \textbf{not} used to claim that the second
study's real fibres have a measured diameter above \(2 eps\); the study
did not measure \(D(z)\) directly. They establish a narrower point: the
empirical certificate was invalid, full coverage at the registered gate
was impossible on the realised excess distribution, and the available
selector did not identify the cases on which abstention would help. The
theory states what must be measured or refined next.

\hypertarget{relation-to-prior-work}{%
\section{8. Relation to prior work}\label{relation-to-prior-work}}

The general mathematical ingredients are donor-owned.
Sufficient-statistic and comparison-of-experiments theory establish that
information is valuable relative to a decision problem rather than in
the abstract. Robust decision theory likewise studies performance when
observations leave sets of compatible states. Interval covering on a
line and the greedy minimum-cover argument are classical algorithmic
facts. We make no novelty claim for those components.

Conformal prediction provides finite-sample marginal coverage under
exchangeability. Recent work on selection-conditional conformal
inference makes the selection event part of the validity target; Jin and
Ren (2025) give exact selection-conditional procedures for broad classes
of selection rules, while Sale and Ramdas (2025) identify failures in
online selective calibration and give exchangeability-preserving
alternatives. Recent conditional-coverage work continues to study how
local reliability can be assessed and compared. Those methods operate at
the statistical calibration layer.

The present paper isolates a complementary deterministic layer:
\textbf{before asking whether calibration is statistically valid, ask
whether the representation identifies the target finely enough for any
fibre-constant certificate at the requested tolerance.} Its residual
contribution is the exact joint calculus linking fibre diameter,
deterministic certificate radius, minimal partition refinement,
separator realizability, and abstention coverage, together with a
preserved empirical failure that motivates the distinction.

\hypertarget{limitations}{%
\section{9. Limitations}\label{limitations}}

\begin{enumerate}
\def\labelenumi{\arabic{enumi}.}
\item
  The core theorems concern finite fibres and a scalar target. Infinite
  spaces require topological or measurable extensions that are not
  supplied here.
\item
  The point-certificate results are deterministic. Randomized
  certificates require an explicit loss and coverage convention.
\item
  The theorem assumes the target values on a fibre when constructing the
  midpoint and computing minimal refinement. Learning or estimating
  those values without leakage is a separate statistical problem.
\item
  Separator realizability says when a declared feature family is
  sufficient. It does not learn the family or price feature acquisition.
\item
  The second-study corpus does not directly measure \(D(z)\) on accepted
  fibres. It is an adverse application record, not empirical
  confirmation of the diameter law.
\item
  The selector-correlation threshold near \texttt{0.85} comes from one
  44-case realised distribution and is not a universal requirement.
\item
  No cross-domain transfer, production benefit, physical quantum
  advantage, computational-hardness result, or broad empirical
  superiority is claimed. The counted empirical studies establish no
  broad empirical superiority or transfer claim.
\item
  Earlier V2 compiler theorems depending on unstated dominance,
  single-block, padding, or cross-gadget assumptions are not part of the
  V3 submission claim. Their historical records are retained for audit.
\end{enumerate}

\hypertarget{discussion-and-conclusion}{%
\section{10. Discussion and
conclusion}\label{discussion-and-conclusion}}

A representation is not simply sufficient or insufficient. Sufficiency
is relative to the question and to the tolerance at which an answer must
be certified. The target diameter of a representation fibre makes that
statement exact. If \(D(z) > 2 eps\), no deterministic certificate that
sees only \texttt{z} can meet error \(eps\); if \(D(z) \le 2 eps\), the
midpoint certificate meets it immediately.

That equivalence turns an impossibility result into a design calculus.
Minimal unconstrained refinement is an exact interval-cover problem.
Realizable refinement is controlled by separator power. Abstention
covers precisely the fibres that are already narrow enough. The
resulting frontier separates three failure modes that are otherwise easy
to mix: a certificate may be badly calibrated, the selector may fail to
identify difficult cases, or the representation may merge
target-separated instances so that no calibration can succeed without
refinement.

The preserved second-study result illustrates why the distinction
matters. Coverage reached every held-out case while conditional validity
failed, and the available selector carried no detected signal about
realised excess. The appropriate response is not to relax the gate after
observing that failure. It is to measure the information problem
prospectively: estimate fibre diameter, test separator sufficiency, or
learn a selector under disjoint custody and then validate the retained
set.

The bounded conclusion is therefore simple: \textbf{certifiability is
governed by within-representation target variation, and the cost of
repairing an insufficient representation can be characterized exactly.}
The external multi-domain question---whether practical learned
representations can approach that information-theoretic
frontier---remains successor science rather than a condition for
submitting the present theory paper.

\hypertarget{tool-use-disclosure}{%
\section{Tool-use disclosure}\label{tool-use-disclosure}}

A generative language model assisted manuscript organization, language
revision, adversarial review, and submission-package preparation.
Scientific claims and final submission decisions remain the
responsibility of the listed author.

\hypertarget{selected-references}{%
\section{Selected references}\label{selected-references}}

\begin{itemize}
\item
  D. Blackwell, \emph{Equivalent Comparisons of Experiments}, Annals of
  Mathematical Statistics 24, 265--272 (1953).
\item
  V. Vovk, A. Gammerman and G. Shafer, \emph{Algorithmic Learning in a
  Random World}, Springer (2005).
\item
  Y. Jin and Z. Ren, \emph{Confidence on the focal: conformal prediction
  with selection-conditional coverage}, Journal of the Royal Statistical
  Society: Series B 87, 1239--1259 (2025). DOI: 10.1093/jrsssb/qkaf016.
\item
  Y. Sale and A. Ramdas, \emph{Online Selective Conformal Prediction:
  Errors and Solutions}, arXiv:2503.16809 (2025).
\item
  Z. Zhou, X. Zhang, C. Tao and Y. Yang, \emph{Conformal Prediction
  Assessment: A Framework for Conditional Coverage Evaluation and
  Selection}, arXiv:2603.27189 (2026).
\end{itemize}

\end{document}
